\chapter{Sprint 1 : Socle Applicatif et Identité Numérique}
\label{chap:sprint1}
\setlength{\fboxsep}{1pt}
\setlength{\fboxrule}{0.4pt}

\section{Introduction}

Avant de pouvoir exploiter les fonctionnalités avancées du \textbf{HR Talent Portal} --- l'intelligence artificielle, les tests gamifiés ou les entretiens techniques --- il est indispensable de poser les fondations sur lesquelles tout le reste reposera. C'est précisément l'objectif de ce premier sprint : construire le socle applicatif et l'identité numérique de la plateforme.

\medskip

Concrètement, ce sprint répond à une question simple : \textit{comment un candidat découvre-t-il les offres de Talan, postule-t-il, et accède-t-il ensuite à son espace personnel en toute sécurité ?} Pour y répondre, nous avons travaillé sur deux axes complémentaires. Le premier, orienté \textbf{vitrine de recrutement}, consiste à déployer le portail public accessible aux visiteurs : une page d'accueil attractive, un catalogue d'offres filtrable par plusieurs critères, et un formulaire de candidature permettant de joindre un CV. Le second axe, orienté \textbf{identité et sécurité}, met en place les mécanismes d'authentification du candidat : connexion par email, connexion rapide par QR~code depuis un mobile, réinitialisation sécurisée du mot de passe, et journalisation de chaque action sensible pour garantir la traçabilité.

\medskip

Dans ce chapitre, nous présenterons d'abord le backlog de ce sprint, puis l'analyse fonctionnelle à travers des cas d'utilisation détaillés et des diagrammes de séquence. Nous exposerons ensuite la conception via le diagramme de classes et les diagrammes d'activité, avant de conclure par la présentation des interfaces réalisées.

\section{Backlog du sprint 1}

Le backlog du sprint~1 (Tableau~\ref{tab:Backlog_sprint1}) regroupe les dix premières User Stories du projet, réparties en quatre domaines fonctionnels : la vitrine publique de recrutement, l'authentification sécurisée, l'espace personnel du candidat, et l'administration des droits d'accès. Chaque User Story est déclinée en tâches concrètes, estimées en jours de développement.
\\
\\
\\
\\

\begin{longtable}{|c|p{5.5cm}|p{6cm}|c|}
\caption{Backlog du sprint 1 : Socle applicatif et identité numérique} \label{tab:Backlog_sprint1} \\
\hline
\rowcolor{gray!20} \textbf{ID} & \textbf{User Story} & \textbf{Tâches techniques} & \textbf{Est. (j)} \\
\hline
\endfirsthead

\multicolumn{4}{c}{\textit{(Suite du backlog du sprint 1)}} \\
\hline
\rowcolor{gray!20} \textbf{ID} & \textbf{User Story} & \textbf{Tâches techniques} & \textbf{Est. (j)} \\
\hline
\endhead

\hline \multicolumn{4}{r}{\textit{Suite sur la page suivante...}} \\
\endfoot

\hline
\endlastfoot

%% ── Portail public ──
\multicolumn{4}{|c|}{\cellcolor{blue!8}\textbf{Portail public --- Vitrine de recrutement}} \\
\hline

\multirow{4}{*}{\textbf{1}} & \multirow{4}{5.5cm}{En tant que visiteur, je souhaite consulter la liste des offres avec des filtres multicritères afin de trouver les postes correspondant à mon profil.} & 1.1 Configuration du site Experience Cloud (domaine, thème LWR, pages publiques). & 1 \\
\cline{3-4}
& & 1.2 Développement du composant LWC \texttt{jobList} (liste d'offres avec pagination). & 1.5 \\
\cline{3-4}
& & 1.3 Développement du composant LWC \texttt{jobFilter} (filtrage par mot-clé, localisation, type de contrat, département). & 1 \\
\cline{3-4}
& & 1.4 Développement du contrôleur Apex \texttt{JobPostingController} (requête SOQL sur Opportunity de Record Type \textit{Job\_Posting}). & 0.5 \\

\hline

\multirow{2}{*}{\textbf{2}} & \multirow{2}{5.5cm}{En tant que visiteur, je souhaite visualiser le détail complet d'une offre afin d'évaluer mon intérêt avant de postuler.} & 2.1 Développement du composant LWC \texttt{jobDetail} (modale avec description, salaire, département). & 1 \\
\cline{3-4}
& & 2.2 Bouton « Postuler » intégré ouvrant le formulaire de candidature. & 0.5 \\

\hline

\multirow{4}{*}{\textbf{3}} & \multirow{4}{5.5cm}{En tant que visiteur, je souhaite soumettre ma candidature et téléverser mon CV (jusqu'à 10~Mo) via un formulaire intuitif.} & 3.1 Développement du composant LWC \texttt{applicationForm} (formulaire multi-champs avec validation). & 1.5 \\
\cline{3-4}
& & 3.2 Développement du contrôleur Apex \texttt{ApplicationController} (création Lead + upload chunké). & 1.5 \\
\cline{3-4}
& & 3.3 Mécanisme d'upload par fragments (chunks de 750~Ko pour respecter la limite Aura de 4~Mo). & 1 \\
\cline{3-4}
& & 3.4 Rattachement du CV via \texttt{ContentDocumentLink} et déclenchement du scoring IA. & 0.5 \\

\hline

%% ── Authentification ──
\multicolumn{4}{|c|}{\cellcolor{green!8}\textbf{Authentification et sécurité}} \\
\hline

\multirow{3}{*}{\textbf{4}} & \multirow{3}{5.5cm}{En tant que candidat, je souhaite me connecter via mon adresse email afin d'accéder à mon espace personnel.} & 4.1 Développement du composant LWC \texttt{loginPage} (modale email/mot de passe). & 1 \\
\cline{3-4}
& & 4.2 Développement de la méthode \texttt{loginWithCredentials()} dans \texttt{QrLoginController} (résolution email~$\rightarrow$~username + \texttt{Site.login()}). & 1 \\
\cline{3-4}
& & 4.3 Configuration du profil \textit{Customer Community Login User} et des Sharing Sets. & 0.5 \\

\hline

\multirow{3}{*}{\textbf{5}} & \multirow{3}{5.5cm}{En tant que candidat, je souhaite réinitialiser mon mot de passe via un lien sécurisé reçu par email.} & 5.1 Développement du composant LWC \texttt{forgetMotDePasse} (saisie email + envoi token). & 1 \\
\cline{3-4}
& & 5.2 Développement du contrôleur Apex \texttt{PasswordResetController} (génération token SHA-256, expiration 24h, usage unique, 3~tentatives max). & 1.5 \\
\cline{3-4}
& & 5.3 Développement du composant LWC \texttt{setPassword} (validation du token + saisie nouveau mot de passe). & 1 \\

\hline

\multirow{2}{*}{\textbf{6}} & \multirow{2}{5.5cm}{En tant que candidat, je souhaite me connecter via un QR~code afin de bénéficier d'un accès rapide depuis mon mobile.} & 6.1 Développement du composant LWC \texttt{qrCodeAuth} (génération + scan QR). & 1 \\
\cline{3-4}
& & 6.2 Développement des méthodes \texttt{generateQrToken()} et \texttt{validateQrToken()} dans \texttt{QrLoginController}. & 1 \\

\hline

%% ── Espace candidat ──
\multicolumn{4}{|c|}{\cellcolor{orange!8}\textbf{Espace candidat B2C}} \\
\hline

\multirow{3}{*}{\textbf{7}} & \multirow{3}{5.5cm}{En tant que candidat, je souhaite visualiser le pipeline de ma candidature en cinq étapes avec mon avancement actuel.} & 7.1 Développement du composant LWC \texttt{candidateDashboard} (pipeline visuel, scores, liens d'évaluation). & 2 \\
\cline{3-4}
& & 7.2 Développement du contrôleur Apex \texttt{CandidateDashboardController} (\texttt{@wire} données candidat). & 1 \\
\cline{3-4}
& & 7.3 Composant \texttt{myApplications} (historique des candidatures). & 1 \\

\hline

\multirow{2}{*}{\textbf{8}} & \multirow{2}{5.5cm}{En tant que candidat, je souhaite consulter et modifier mes informations personnelles.} & 8.1 Onglet « Profil » dans \texttt{candidateDashboard} (mode lecture/édition). & 1 \\
\cline{3-4}
& & 8.2 Mise à jour des champs Contact via méthode Apex dédiée. & 0.5 \\

\hline

%% ── Administration & Audit ──
\multicolumn{4}{|c|}{\cellcolor{purple!8}\textbf{Audit et administration}} \\
\hline

\multirow{3}{*}{\textbf{9}} & \multirow{3}{5.5cm}{En tant que responsable RH, je souhaite que chaque opération sensible soit journalisée dans un audit trail.} & 9.1 Création de l'objet \texttt{Security\_Audit\_Log\_\_c} (Action, Email, Severity, Timestamp, Details). & 0.5 \\
\cline{3-4}
& & 9.2 Instrumentation des contrôleurs d'authentification (login, reset, QR) pour journaliser chaque action. & 1 \\
\cline{3-4}
& & 9.3 Niveaux de sévérité : INFO, WARNING, CRITICAL. & 0.5 \\

\hline

\multirow{3}{*}{\textbf{10}} & \multirow{3}{5.5cm}{En tant que responsable RH, je souhaite gérer les habilitations utilisateurs via des PermissionSets dédiés.} & 10.1 Création du PermissionSet \texttt{Candidate\_Portal\_Access} (FLS candidat B2C). & 0.5 \\
\cline{3-4}
& & 10.2 Création du PermissionSet \texttt{Security\_Objects\_Full\_Access} (objets de sécurité). & 0.5 \\
\cline{3-4}
& & 10.3 Configuration de l'accès Guest User aux classes Apex exposées (\texttt{@AuraEnabled}). & 0.5 \\

\hline
\end{longtable}

\section{Analyse fonctionnelle du sprint 1}

Après avoir identifié les fonctionnalités à développer dans ce sprint, nous les formalisons à travers des cas d'utilisation détaillés. L'objectif est de décrire, de manière claire et structurée, les interactions entre chaque acteur et le système, afin de guider ensuite la conception et le développement. Nous présentons d'abord le diagramme de cas d'utilisation raffiné, puis les descriptions textuelles des cas principaux, et enfin les diagrammes de séquence qui illustrent la chronologie des échanges.

\subsection{Diagramme de cas d'utilisation raffiné du sprint 1}

Le diagramme de cas d'utilisation du sprint~1 (figure~\ref{fig:uc_sprint1}) fait intervenir quatre acteurs. Le \textbf{Visiteur} navigue sur le portail public pour consulter les offres et déposer sa candidature. Le \textbf{Candidat}, une fois authentifié, accède à son espace personnel pour suivre l'avancement de son dossier et gérer son profil. Le \textbf{Responsable RH} administre les droits d'accès et les habilitations. Enfin, le \textbf{Système} assure en arrière-plan la journalisation automatique de toute action sensible (connexion, réinitialisation de mot de passe, téléversement de fichier).

\medskip

Les relations d'inclusion ($\ll$include$\gg$) expriment une règle métier simple : pour accéder à son espace personnel, le candidat doit obligatoirement s'être authentifié au préalable.

\begin{figure}[hbt!]
  \centering
 
   \fbox{\includegraphics[width=14cm]{images/usecaseSprint1.png}}
  \caption{Diagramme de cas d'utilisation raffiné du sprint 1}
  \label{fig:uc_sprint1}
\end{figure}
\FloatBarrier

\subsection{Cas d'utilisation « Consulter les offres d'emploi »}

\subsubsection{Description textuelle du cas d'utilisation « Consulter les offres d'emploi »}
(Voir Tableau~\ref{tab:cu_consulter_offres})

\begin{xltabular}{\linewidth}{|X|}
    \hline
    \textbf{Cas d'utilisation} : Consulter les offres d'emploi \\
    \hline
    \textbf{Acteur} : Visiteur \\
    \hline
    \textbf{Précondition} : Le visiteur accède au portail public via son navigateur. \\
    \hline
    \textbf{Postcondition} : La liste des offres correspondant aux critères est affichée. \\
    \hline
    \textbf{Scénario nominal} : \\
    1. Le visiteur accède à la page d'accueil du portail HR Talent Portal. \\
    2. Le système affiche la liste des offres d'emploi ouvertes, triées par date de publication décroissante (les plus récentes en premier). \\
    3. Le visiteur applique un ou plusieurs filtres : mot-clé, localisation, type de contrat ou département. \\
    4. Le système met à jour les résultats en temps réel, sans rechargement de la page. \\
    5. Le visiteur sélectionne une offre pour en consulter le détail. \\
    6. Le système affiche une fiche détaillée présentant le titre, la description complète, le salaire proposé, le département et la localisation. \\
    \hline
    \textbf{Scénario alternatif} : \\
    3.a Aucune offre ne correspond aux critères de filtrage. \\
    \hspace{0.5cm} 3.a.1 Le système affiche un message indiquant qu'aucun résultat n'a été trouvé. \\
    \hline
\caption{Description textuelle du CU « Consulter les offres d'emploi »}
\label{tab:cu_consulter_offres}
\end{xltabular}

\subsection{Cas d'utilisation « Soumettre une candidature »}

\subsubsection{Description textuelle du cas d'utilisation « Soumettre une candidature »}
(Voir Tableau~\ref{tab:cu_soumettre_candidature})

\begin{xltabular}{\linewidth}{|X|}
    \hline
    \textbf{Cas d'utilisation} : Soumettre une candidature \\
    \hline
    \textbf{Acteur} : Visiteur \\
    \hline
    \textbf{Précondition} : Le visiteur a consulté le détail d'une offre d'emploi et cliqué sur « Postuler ». \\
    \hline
    \textbf{Postcondition} : Un dossier de candidature est créé dans le système avec le CV rattaché. L'analyse automatique du CV par l'IA est déclenchée en arrière-plan. \\
    \hline
    \textbf{Scénario nominal} : \\
    1. Le système affiche le formulaire de candidature, pré-rempli avec la référence de l'offre concernée. \\
    2. Le visiteur renseigne ses informations personnelles : nom, prénom, email, téléphone, adresse, code postal et un message de motivation. \\
    3. Le visiteur joint son CV (formats acceptés : PDF, DOC, DOCX, TXT ; taille maximale : 10~Mo). \\
    4. Le visiteur accepte les conditions générales et soumet le formulaire. \\
    5. Le système crée le dossier de candidature dans la base de données. \\
    6. Si le fichier est volumineux (supérieur à 1~Mo), le système procède à un envoi par morceaux afin de contourner les limites techniques de la plateforme. Sinon, le fichier est téléversé directement. \\
    7. Le système rattache le document au dossier du candidat. \\
    8. Le système déclenche l'analyse IA du CV en arrière-plan (traitement asynchrone). \\
    9. Le système affiche un message de confirmation au visiteur. \\
    \hline
    \textbf{Scénario alternatif} : \\
    4.a Le visiteur n'a pas rempli tous les champs obligatoires ou l'email est invalide. \\
    \hspace{0.5cm} 4.a.1 Le système affiche les messages d'erreur sous les champs concernés. \\
    3.a Le fichier dépasse 10~Mo ou son format n'est pas accepté. \\
    \hspace{0.5cm} 3.a.1 Le système affiche un message d'erreur spécifiant les formats et la taille autorisés. \\
    \hline
\caption{Description textuelle du CU « Soumettre une candidature »}
\label{tab:cu_soumettre_candidature}
\end{xltabular}

\subsection{Cas d'utilisation « S'authentifier »}

\subsubsection{Description textuelle du cas d'utilisation « S'authentifier par email »}
(Voir Tableau~\ref{tab:cu_authentifier})

\begin{xltabular}{\linewidth}{|X|}
    \hline
    \textbf{Cas d'utilisation} : S'authentifier par email \\
    \hline
    \textbf{Acteur} : Candidat \\
    \hline
    \textbf{Précondition} : Le candidat dispose d'un compte portail (créé automatiquement lorsque son CV a été qualifié par l'IA). \\
    \hline
    \textbf{Postcondition} : Le candidat est authentifié et accède à son tableau de bord personnel. \\
    \hline
    \textbf{Scénario nominal} : \\
    1. Le candidat clique sur \og Se connecter \fg dans le bandeau de navigation du portail. \\
    2. Le système affiche la fenêtre de connexion. \\
    3. Le candidat saisit son adresse email et son mot de passe. \\
    4. Le système recherche le compte correspondant à l'adresse email fournie. Cette résolution est nécessaire car Salesforce utilise en interne un identifiant technique différent de l'email --- le candidat, lui, n'a besoin de connaître que son email. \\
    5. Le système vérifie les identifiants auprès du moteur d'authentification de Salesforce. \\
    6. Le candidat est redirigé vers son tableau de bord personnel. \\
    7. Le système enregistre cette connexion dans le journal d'audit (niveau informationnel). \\
    \hline
    \textbf{Scénario alternatif} : \\
    5.a L'email ou le mot de passe est incorrect. \\
    \hspace{0.5cm} 5.a.1 Le système affiche un message générique : \og Email ou mot de passe incorrect \fg. Ce message est volontairement identique quel que soit le motif d'échec, afin d'empêcher un attaquant de déduire l'existence d'un compte (principe d'anti-énumération). \\
    \hspace{0.5cm} 5.a.2 Le système enregistre la tentative échouée dans le journal d'audit (niveau avertissement). \\
    \hline
\caption{Description textuelle du CU « S'authentifier par email »}
\label{tab:cu_authentifier}
\end{xltabular}

\subsubsection{Description textuelle du cas d'utilisation « S'authentifier par QR~code »}
(Voir Tableau~\ref{tab:cu_qr_auth})

\begin{xltabular}{\linewidth}{|X|}
    \hline
    \textbf{Cas d'utilisation} : S'authentifier par QR code \\
    \hline
    \textbf{Acteur} : Candidat \\
    \hline
    \textbf{Précondition} : Le candidat dispose d'un compte portail actif et accède au portail depuis un terminal mobile ou un second écran. \\
    \hline
    \textbf{Postcondition} : Le candidat est authentifié sans saisie manuelle d'identifiants. \\
    \hline
    \textbf{Scénario nominal} : \\
    1. Le candidat sélectionne l'option \og Connexion par QR~Code \fg sur la page de login. \\
    2. Le système génère un jeton unique et sécurisé, valable pendant 5~minutes. Ce jeton est stocké sous forme chiffrée (hash SHA-256) dans la base de données. \\
    3. Le portail affiche un QR~code encodant ce jeton. \\
    4. Le candidat scanne le QR~code avec son téléphone mobile. \\
    5. L'application mobile valide le jeton auprès du serveur. \\
    6. Le navigateur, qui vérifie périodiquement l'état du jeton, détecte la validation et ouvre automatiquement la session du candidat. \\
    \hline
    \textbf{Scénario alternatif} : \\
    6.a Le jeton expire avant le scan (délai de 5~minutes dépassé). \\
    \hspace{0.5cm} 6.a.1 Le système affiche un message invitant le candidat à générer un nouveau QR~code. \\
    \hline
\caption{Description textuelle du CU « S'authentifier par QR code »}
\label{tab:cu_qr_auth}
\end{xltabular}

\subsection{Cas d'utilisation « Réinitialiser le mot de passe »}

\subsubsection{Description textuelle du cas d'utilisation « Réinitialiser le mot de passe »}
(Voir Tableau~\ref{tab:cu_reset_password})

\begin{xltabular}{\linewidth}{|X|}
    \hline
    \textbf{Cas d'utilisation} : Réinitialiser le mot de passe \\
    \hline
    \textbf{Acteur} : Candidat \\
    \hline
    \textbf{Précondition} : Le candidat possède un compte portail mais a oublié son mot de passe. \\
    \hline
    \textbf{Postcondition} : Le mot de passe est mis à jour et le candidat peut s'authentifier. \\
    \hline
    \textbf{Scénario nominal} : \\
    1. Le candidat clique sur \og Mot de passe oublié~? \fg depuis la fenêtre de connexion. \\
    2. Le système affiche un formulaire demandant l'adresse email. \\
    3. Le candidat saisit son email et soumet la demande. \\
    4. Le système génère un jeton aléatoire de haute entropie (256~bits), calcule son empreinte numérique (hash SHA-256) et enregistre les métadonnées associées : empreinte du jeton, email, délai d'expiration de 24~heures, compteur de tentatives à zéro et indicateur d'usage unique. Seul le hash est conservé dans la base de données --- le jeton en clair n'est jamais persisté, ce qui protège les utilisateurs même en cas de compromission de la base. \\
    5. Le système envoie un email contenant un lien sécurisé avec le jeton (usage unique). \\
    6. Le système enregistre cette opération dans le journal d'audit. \\
    7. Le candidat clique sur le lien reçu par email. \\
    8. Le système affiche la page de définition du nouveau mot de passe et vérifie la validité du jeton en contrôlant quatre conditions : (a)~l'empreinte correspond, (b)~le délai de 24~heures n'est pas dépassé, (c)~le jeton n'a pas déjà été utilisé, et (d)~le nombre de tentatives est inférieur ou égal à 3. \\
    9. Le candidat saisit son nouveau mot de passe, qui doit respecter des critères de robustesse : au moins 8~caractères, une majuscule, une minuscule, un chiffre et un caractère spécial. \\
    10. Le système vérifie la robustesse, applique le changement et marque le jeton comme utilisé pour empêcher toute réutilisation. \\
    11. Le système enregistre le succès de la réinitialisation dans le journal d'audit. \\
    \hline
    \textbf{Scénario alternatif} : \\
    8.a Le jeton est invalide, expiré ou déjà utilisé. \\
    \hspace{0.5cm} 8.a.1 Le système affiche un message générique : \og Lien invalide ou expiré \fg. Comme pour la connexion, le message est identique quel que soit le motif réel, afin de ne pas donner d'indice à un éventuel attaquant. \\
    \hspace{0.5cm} 8.a.2 Le système enregistre la tentative dans le journal d'audit (niveau avertissement). \\
    8.b Le nombre de tentatives dépasse 3. \\
    \hspace{0.5cm} 8.b.1 Le système invalide définitivement le jeton et enregistre l'événement (niveau critique). \\
    9.a Le mot de passe ne respecte pas les critères de robustesse. \\
    \hspace{0.5cm} 9.a.1 Le système affiche les critères non satisfaits. \\
    \hline
\caption{Description textuelle du CU « Réinitialiser le mot de passe »}
\label{tab:cu_reset_password}
\end{xltabular}

\subsection{Cas d'utilisation « Gérer les habilitations »}

\subsubsection{Description textuelle du cas d'utilisation « Gérer les habilitations »}
(Voir Tableau~\ref{tab:cu_gerer_habilitations})

\begin{xltabular}{\linewidth}{|X|}
    \hline
    \textbf{Cas d'utilisation} : Gérer les habilitations \\
    \hline
    \textbf{Acteur} : Responsable RH \\
    \hline
    \textbf{Précondition} : Le responsable RH est authentifié dans l'interface d'administration Salesforce. \\
    \hline
    \textbf{Postcondition} : Les droits d'accès sont correctement attribués ; les visiteurs non authentifiés ne peuvent accéder qu'aux fonctionnalités publiques. \\
    \hline
    \textbf{Scénario nominal} : \\
    1. Le responsable RH accède à la page de gestion des droits d'accès dans la configuration Salesforce. \\
    2. Le système affiche les jeux de permissions existants, organisés par profil d'utilisateur. \\
    3. Le responsable RH attribue le jeu de permissions \og Accès Portail Candidat \fg aux utilisateurs du portail B2C. Ce jeu définit précisément quels champs et quels objets un candidat peut consulter ou modifier depuis son espace personnel. \\
    4. Le responsable RH configure les accès du visiteur non authentifié : celui-ci ne peut interagir qu'avec les fonctionnalités publiques du portail (consultation des offres, soumission de candidature). \\
    5. Le système enregistre les modifications. \\
    \hline
\caption{Description textuelle du CU « Gérer les habilitations »}
\label{tab:cu_gerer_habilitations}
\end{xltabular}

\subsection{Diagrammes de séquence système du sprint 1}

Les diagrammes de séquence permettent de visualiser la chronologie des échanges entre l'utilisateur et le système pour chaque fonctionnalité clé. Ils montrent \textit{qui} communique avec \textit{qui}, et dans quel ordre, ce qui facilite la compréhension du comportement attendu avant de passer au développement.

\subsubsection{Diagramme de séquence \og Soumettre une candidature \fg}

La figure~\ref{fig:seq_candidature} retrace le parcours d'un visiteur qui soumet sa candidature. Le flux commence par le remplissage du formulaire, se poursuit avec la création du dossier candidat dans la base de données, puis le téléversement du CV --- découpé en morceaux si le fichier est volumineux --- et se conclut par le déclenchement automatique de l'analyse IA en arrière-plan. Ce dernier point est essentiel : dès qu'un CV est déposé, le système lance immédiatement son évaluation, sans aucune intervention humaine.

\begin{figure}[hbt!]
  \centering
 
   \fbox{\includegraphics[width=16cm]{images/diagSeqsoumettrecond.png}}
  \caption{Diagramme de séquence système du CU « Soumettre une candidature »}
  \label{fig:seq_candidature}
\end{figure}
\FloatBarrier

\subsubsection{Diagramme de séquence \og S'authentifier \fg}

La figure~\ref{fig:seq_authentifier} illustre le processus de connexion par email. Le candidat fournit son adresse email et son mot de passe. Le système traduit l'email en identifiant technique Salesforce --- une étape transparente pour l'utilisateur mais nécessaire car Salesforce distingue l'email de l'identifiant de connexion --- puis vérifie les identifiants. Chaque tentative, qu'elle aboutisse ou échoue, est consignée dans le journal d'audit pour assurer la traçabilité complète des accès.

\begin{figure}[hbt!]
  \centering
 
   \fbox{\includegraphics[width=16cm]{images/diagSeqAuthentifier.png}}
  \caption{Diagramme de séquence système du CU « S'authentifier par email »}
  \label{fig:seq_authentifier}
\end{figure}
\FloatBarrier

\subsubsection{Diagramme de séquence \og Réinitialiser le mot de passe \fg}

La figure~\ref{fig:seq_reset_password} retrace le flux complet de réinitialisation, qui se déroule en deux temps bien distincts. Dans un premier temps, le candidat demande la réinitialisation : le système génère un jeton sécurisé et l'envoie par email. Dans un second temps, le candidat clique sur le lien reçu : le système vérifie la validité du jeton (empreinte, expiration, usage unique, nombre de tentatives), puis autorise la définition d'un nouveau mot de passe. Ce mécanisme en deux étapes garantit que seule la personne ayant accès à la boîte email peut modifier le mot de passe.

\begin{figure}[hbt!]
  \centering

   \fbox{\includegraphics[width=14cm]{images/diagSeqReset.png}}
  \caption{Diagramme de séquence système du CU « Réinitialiser le mot de passe »}
  \label{fig:seq_reset_password}
\end{figure}
\FloatBarrier

\section{Conception du sprint 1}

Cette section présente la modélisation structurelle et comportementale des fonctionnalités développées dans ce sprint, à travers un diagramme de classes et des diagrammes d'activité.

\subsection{Diagramme de classes du sprint 1}

Le diagramme de classes du sprint~1 (figure~\ref{fig:diagclass_sprint1}) modélise les entités de données nécessaires au socle applicatif et à l'identité numérique. Il fait intervenir neuf entités, que l'on peut regrouper en trois familles selon leur rôle dans le système.

\medskip

\textbf{La famille « candidature ».} Tout commence lorsqu'un visiteur dépose sa candidature : le système crée un \textbf{Lead}, c'est-à-dire un dossier initial contenant les informations du candidat (nom, email, téléphone) ainsi qu'une référence vers l'offre concernée. Le CV du candidat est stocké dans l'entité \textbf{ContentVersion} (qui gère les versions successives de tout document dans Salesforce) et rattaché au dossier via l'entité de liaison \textbf{ContentDocumentLink}. Quant aux offres d'emploi, elles sont portées par l'entité \textbf{Opportunity} --- plus précisément par une variante spécialisée appelée \textit{Job\_Posting}, obtenue grâce au mécanisme de \textbf{Record Types}, qui permet de réutiliser un même objet Salesforce pour des usages différents.

\medskip

\textbf{La famille « identité ».} Lorsque le candidat est qualifié par l'IA (cette conversion sera traitée au Sprint~2), un \textbf{Contact} et un \textbf{Account} sont créés automatiquement à partir du Lead. L'entité \textbf{User} représente ensuite le compte de connexion au portail, relié au Contact. C'est cette entité qui permet au candidat de s'authentifier et d'accéder à son espace personnel.

\medskip

\textbf{La famille « sécurité ».} Trois entités dédiées complètent le modèle. L'entité de gestion des jetons de réinitialisation stocke l'empreinte SHA-256 de chaque jeton, sa date d'expiration, son indicateur d'usage unique et le compteur de tentatives --- tous les garde-fous nécessaires pour protéger le processus de récupération de mot de passe. L'entité d'authentification par QR~code gère les jetons éphémères de 5~minutes utilisés pour la connexion alternative depuis un mobile. Enfin, le journal d'audit enregistre chaque action sensible (connexion, réinitialisation, échec) avec un horodatage, le niveau de sévérité et les détails contextuels, assurant ainsi la traçabilité complète des opérations de sécurité.
\begin{figure}[hbt!]
  \centering
  
   \fbox{\includegraphics[width=17cm]{images/diagClassesprint1.png}}
  \caption{Diagramme de classes du sprint 1}
  \label{fig:diagclass_sprint1}
\end{figure}
\FloatBarrier

\subsection{Diagrammes d'activité du sprint 1}

\subsubsection{Diagramme d'activité « Soumettre une candidature avec upload de CV »}

Le processus de soumission d'une candidature (figure~\ref{fig:act_candidature}) illustre un défi technique concret : comment permettre à un visiteur de téléverser un CV pouvant atteindre 10~Mo, alors que la plateforme Salesforce impose une limite de 4~Mo par requête~? La solution adoptée repose sur un découpage intelligent du fichier en morceaux. Le flux se décompose en trois phases :

\begin{enumerate}
    \item \textbf{Validation et création du dossier :} le formulaire vérifie d'abord que tous les champs obligatoires sont remplis et que l'email est valide. Si tout est correct, le système crée le dossier candidat dans la base de données.
    \item \textbf{Téléversement du CV :} si le fichier dépasse 1~Mo, il est automatiquement découpé en fragments de 750~Ko. Chaque fragment est envoyé séparément et reconstitué côté serveur. Une barre de progression informe le candidat de l'avancement en temps réel. Pour les fichiers plus petits, l'envoi est direct et instantané.
    \item \textbf{Déclenchement de l'analyse IA :} une fois le CV rattaché au dossier, le système programme automatiquement l'analyse par intelligence artificielle. Ce traitement s'exécute en arrière-plan, sans bloquer l'interface du visiteur qui reçoit immédiatement un message de confirmation.
\end{enumerate}

\begin{figure}[hbt!]
  \centering
  
  \fbox{\includegraphics[width=15cm]{images/activitéSoumettreCondi.png}}
  \caption{Diagramme d'activité « Soumettre une candidature avec upload chunké »}
  \label{fig:act_candidature}
\end{figure}
\FloatBarrier

\subsubsection{Diagramme d'activité « Réinitialiser le mot de passe »}

Le processus de réinitialisation (figure~\ref{fig:act_reset}) met en œuvre une série de garde-fous de sécurité inspirés des recommandations OWASP, le standard international de sécurité des applications web. L'idée centrale est de ne jamais faire confiance au lien seul : il faut vérifier plusieurs conditions indépendantes avant d'autoriser le changement de mot de passe.

\begin{enumerate}
    \item \textbf{Phase de demande :} le candidat fournit son adresse email. Le système génère un jeton aléatoire, en conserve uniquement l'empreinte numérique (hash SHA-256) --- jamais le jeton en clair --- et envoie un email contenant le lien de réinitialisation. Même si l'email n'existe pas dans le système, le même message de confirmation est affiché, afin qu'un attaquant ne puisse pas deviner quels comptes existent.
    \item \textbf{Phase de validation :} lorsque le candidat clique sur le lien, le système effectue quatre vérifications : (1)~l'empreinte du jeton correspond, (2)~le délai de 24~heures n'est pas dépassé, (3)~le jeton n'a pas déjà été utilisé, (4)~le nombre de tentatives ne dépasse pas~3. Si toutes ces conditions sont réunies, le candidat peut définir un nouveau mot de passe qui doit satisfaire les critères de robustesse.
\end{enumerate}

Chaque opération, qu'elle réussisse ou échoue, est systématiquement tracée dans le journal d'audit avec le niveau de sévérité adéquat (information pour les succès, avertissement pour les échecs, critique en cas de dépassement du seuil de tentatives).

\begin{figure}[hbt!]
  \centering
  % TODO : Insérer le diagramme d'activité « Réinitialiser le mot de passe »
   \fbox{\includegraphics[width=15cm]{images/activitéRest.png}}
  \caption{Diagramme d'activité « Réinitialiser le mot de passe »}
  \label{fig:act_reset}
\end{figure}
\FloatBarrier

\section{Réalisation technique}

Cette section présente les livrables concrets du sprint~1 : les composants développés, l'organisation de l'architecture et les interfaces utilisateur telles qu'elles apparaissent aux différents acteurs.

\subsection{Architecture technique des composants}

Le sprint~1 a produit \textbf{14~composants d'interface} (Lightning Web Components), que l'on peut regrouper en quatre familles selon leur fonction :

\begin{itemize}
    \item \textbf{Composants de structure} (4~composants) : la barre de navigation avec le logo Talan et le bouton de connexion, le pied de page, la bannière d'accueil avec un appel à l'action invitant à découvrir les offres, et la barre de recherche globale.
    
    \item \textbf{Composants métier --- offres et candidature} (4~composants) : la grille d'affichage des offres d'emploi, le panneau de filtres multicritères (mot-clé, localisation, contrat, département), la fiche détaillée d'une offre avec le bouton « Postuler », et le formulaire de candidature avec téléversement progressif du CV.
    
    \item \textbf{Composants d'identité} (4~composants) : la fenêtre de connexion par email et mot de passe, la page de demande de réinitialisation, la page de définition du nouveau mot de passe, et le module de connexion par QR~code.
    
    \item \textbf{Composants d'espace personnel} (2~composants) : le tableau de bord du candidat (avec onglets Vue~d'ensemble, Candidatures et Profil) et la liste de l'historique des candidatures.
\end{itemize}

\medskip

Côté serveur, cinq contrôleurs assurent la logique métier : un pour l'exposition des offres en lecture seule, un pour la création des candidatures et la gestion du téléversement, un pour l'authentification (par email et par QR~code), un pour la réinitialisation sécurisée du mot de passe, et un pour l'envoi des emails de bienvenue.

\subsection{Interface « Page d'accueil du portail »}

La page d'accueil du portail HR Talent Portal constitue le premier point de contact entre Talan et les talents externes. Construite sur le moteur \textbf{LWR (Lightning Web Runtime)} pour sa rapidité de chargement, elle se compose de trois zones :

\begin{enumerate}
    \item \textbf{Le header :} barre de navigation responsive intégrant le logo Talan, les liens de navigation et un bouton « Se connecter » qui déclenche l'événement global \texttt{open-login-modal}.
    \item \textbf{La bannière hero :} composant \texttt{startupHero} présentant un message d'accroche avec un appel à l'action dirigeant vers la section des offres.
    \item \textbf{Le catalogue d'offres :} composants \texttt{jobFilter} (panneau de filtres) et \texttt{jobList} (grille d'offres) fonctionnant en tandem. Le filtre émet un événement personnalisé \texttt{filter} contenant les critères (keyword, localisation, contrat, département) ; la liste réagit en filtrant les résultats en temps réel côté client, sans appel serveur supplémentaire.
\end{enumerate}

\begin{figure}[hbt!]
  \centering
  % TODO : Insérer capture de la page d'accueil
   \fbox{\includegraphics[width=\dimexpr\textwidth-4pt\relax]{images/Capture d'écran 2026-04-16 150916.png}}
  \caption{Interface de la page d'accueil du portail HR Talent Portal}
  \label{fig:accueil_portail}
\end{figure}
\FloatBarrier

\subsection{Interface « Consulter les offres d'emploi »}

Les offres d'emploi sont présentées sous forme de cartes visuelles, chacune affichant les informations essentielles : titre du poste, département, localisation, type de contrat et un aperçu de la description. Lorsque le visiteur souhaite en savoir plus, un bouton « Voir plus » ouvre une fiche détaillée (figure~\ref{fig:detail_offre}) présentant la description complète de l'offre, le salaire proposé et un bouton « Postuler maintenant » qui ouvre directement le formulaire de candidature.

\medskip

Quatre filtres combinables permettent d'affiner la recherche :
\begin{itemize}
    \item \textbf{Recherche textuelle :} filtrage par mot-clé sur le nom et la description du poste.
    \item \textbf{Localisation :} sélection parmi les villes où Talan est présent.
    \item \textbf{Type de contrat :} CDI, CDD, Stage, Freelance ou Alternance.
    \item \textbf{Département :} les différents pôles de Talan (IT, Consulting, Finance, etc.).
\end{itemize}
L'ensemble des filtres fonctionne en temps réel : les résultats se mettent à jour instantanément au fur et à mesure de la saisie, offrant une expérience de recherche fluide.

\begin{figure}[htp]
\begin{minipage}[b]{0.5\linewidth}
\centering

 \fbox{\includegraphics[width=\dimexpr\textwidth-4pt\relax]{images/offres.png}}
\caption{Interface de la liste des offres d'emploi}
\label{fig:liste_offres}
\end{minipage}
\hspace{0.25cm}
\begin{minipage}[b]{0.5\linewidth}
\centering

 \fbox{\includegraphics[width=\dimexpr\textwidth-4pt\relax]{images/detailleoffres.png}}
\caption{Interface du détail d'une offre d'emploi}
\label{fig:detail_offre}
\end{minipage}
\end{figure}
\FloatBarrier

\subsection{Interface « Soumettre une candidature »}

Le formulaire de candidature (figure~\ref{fig:form_candidature}) se présente sous forme de fenêtre accessible depuis la fiche détaillée d'une offre. Le candidat y renseigne ses informations : prénom, nom, adresse email, téléphone, adresse postale, code postal et un message de motivation optionnel. Un sélecteur de fichier permet de joindre le CV, avec une validation automatique du format (PDF, DOC, DOCX, TXT) et de la taille (10~Mo maximum).

\medskip

Lorsque le fichier est volumineux, une barre de progression dynamique informe le candidat de l'avancement du téléversement en temps réel. La soumission est conditionnée par l'acceptation des conditions générales. Un indicateur de chargement visuel est affiché pendant le traitement, et un message de confirmation apparaît une fois la candidature envoyée avec succès (figure~\ref{fig:upload_progress}).

\begin{figure}[htp]
\begin{minipage}[b]{0.5\linewidth}
\centering

 \fbox{\includegraphics[width=\dimexpr\textwidth-4pt\relax]{images/soumettreCondidature.png}}
\caption{Interface du formulaire de candidature}
\label{fig:form_candidature}
\end{minipage}
\hspace{0.25cm}
\begin{minipage}[b]{0.5\linewidth}
\centering

 \fbox{\includegraphics[width=\dimexpr\textwidth-4pt\relax]{images/condidatureenvoyé.png}}
\caption{Condidature Envoyée}
\label{fig:upload_progress}
\end{minipage}
\end{figure}
\FloatBarrier

\subsection{Interface « S'authentifier »}

La page de connexion (figure~\ref{fig:login_page}) offre au candidat deux manières de s'identifier. La première, classique, consiste à saisir son adresse email et son mot de passe. Trois actions sont proposées : se connecter, réinitialiser son mot de passe en cas d'oubli, ou opter pour une connexion par QR~code. Cette fenêtre est accessible depuis n'importe quelle page du portail via le bouton \og Se connecter \fg du bandeau de navigation.

\medskip

La seconde option (figure~\ref{fig:qr_login}) s'adresse aux candidats souhaitant une connexion rapide depuis un téléphone mobile. Le portail affiche un QR~code unique, valable 5~minutes. Le candidat le scanne avec son téléphone, ce qui valide automatiquement la session sur le navigateur --- sans aucune saisie manuelle d'identifiants.

\begin{figure}[htp]
\begin{minipage}[b]{0.5\linewidth}
\centering
 \fbox{\includegraphics[width=\dimexpr\textwidth-4pt\relax]{images/authentification.png}}
\caption{Interface de connexion (email / mot de passe)}
\label{fig:login_page}
\end{minipage}
\hspace{0.25cm}
\begin{minipage}[b]{0.5\linewidth}
\centering

 \fbox{\includegraphics[width=\dimexpr\textwidth-4pt\relax]{images/QRCode.png}}
\caption{Interface de connexion par QR code}
\label{fig:qr_login}
\end{minipage}
\end{figure}
\FloatBarrier

\subsection{Interface « Réinitialiser le mot de passe »}

Le processus de réinitialisation se déroule en deux écrans successifs. Le premier (figure~\ref{fig:forget_password}) invite le candidat à saisir son adresse email. Après soumission, un message de confirmation est affiché \textit{dans tous les cas}, même si l'adresse n'est pas reconnue : cette approche empêche un attaquant de deviner quels comptes existent dans le système.

\medskip

Le second écran (figure~\ref{fig:set_password}), accessible via le lien reçu par email, présente deux champs (nouveau mot de passe et confirmation). Une validation en temps réel guide le candidat en affichant, au fur et à mesure de la saisie, quels critères de robustesse sont satisfaits : longueur minimale de 8~caractères, présence d'une majuscule, d'une minuscule, d'un chiffre et d'un caractère spécial.

\begin{figure}[htp]
\begin{minipage}[b]{0.5\linewidth}
\centering
% TODO : Insérer capture de la page « Mot de passe oublié »
 \fbox{\includegraphics[width=\dimexpr\textwidth-4pt\relax]{images/reset-mot-depasse.png}}
\caption{Interface « Mot de passe oublié »}
\label{fig:forget_password}
\end{minipage}
\hspace{0.25cm}
\begin{minipage}[b]{0.5\linewidth}
\centering

 \fbox{\includegraphics[width=\dimexpr\textwidth-4pt\relax]{images/RedefinerMDP.png}}
\caption{Interface de définition du nouveau mot de passe}
\label{fig:set_password}
\end{minipage}
\end{figure}
\FloatBarrier

\subsection{Interface « Tableau de bord candidat »}

Le tableau de bord (figure~\ref{fig:dashboard_candidat}) constitue l'espace personnel du candidat une fois authentifié. Il est organisé en trois onglets, chacun répondant à un besoin différent :

\begin{enumerate}
    \item \textbf{Vue d'ensemble :} c'est la page principale. Elle affiche le pipeline de recrutement sous forme d'une frise horizontale en cinq étapes : Screening~IA, Gaming, Technique, Architecture et Décision~RH. Chaque étape est visuellement distinguée par un code couleur : gris pour les étapes à venir, bleu pour l'étape en cours et vert pour les étapes complétées. Cette représentation visuelle permet au candidat de savoir instantanément où il en est dans le processus. Les scores obtenus et les liens vers les évaluations planifiées sont affichés sous la frise.
    
    \item \textbf{Candidatures :} cet onglet liste l'ensemble des candidatures passées et en cours, avec pour chacune le statut actuel, la date de soumission, le nom de l'offre et un badge coloré indiquant le score de matching (rouge pour les scores faibles, orange pour les scores intermédiaires, vert pour les scores qualifiants).
    
    \item \textbf{Profil :} le candidat peut consulter et modifier ses informations personnelles (nom, prénom, email, téléphone). Un bouton « Modifier » active le mode édition, et les changements sont enregistrés immédiatement.
\end{enumerate}

Si un visiteur non authentifié tente d'accéder à cette page, le système le redirige automatiquement vers la page de connexion.

\begin{figure}[hbt!]
  \centering
  
  \fbox{\includegraphics[width=\dimexpr\textwidth-4pt\relax]{images/condidatdasboard.png}}
  \caption{Interface du tableau de bord candidat (vue d'ensemble)}
  \label{fig:dashboard_candidat}
\end{figure}
\FloatBarrier

\subsection{Mécanismes de sécurité implémentés}

La sécurité n'est pas un ajout tardif mais un axe structurant dès ce premier sprint. En effet, un portail de recrutement manipule des données personnelles sensibles (CV, coordonnées, parcours professionnel) et doit inspirer confiance aux candidats. Le tableau~\ref{tab:securite_sprint1} synthétise les huit mécanismes de protection mis en \oe{}uvre, chacun répondant à une menace concrète.

\begin{longtable}{|p{4cm}|p{4cm}|p{6.5cm}|}
\caption{Mécanismes de sécurité implémentés au sprint 1} \label{tab:securite_sprint1} \\
\hline
\rowcolor{gray!20} \textbf{Mécanisme} & \textbf{Périmètre} & \textbf{Description} \\
\hline
\endfirsthead
\endhead

Token SHA-256 à usage unique & Réinitialisation du mot de passe & Le système ne conserve jamais le jeton en clair : seule son empreinte numérique (SHA-256) est stockée. Chaque jeton expire après 24~heures et ne peut être utilisé qu'une seule fois, avec un maximum de 3~tentatives. \\
\hline

Anti-énumération & Connexion et réinitialisation & Les messages d'erreur sont volontairement identiques quel que soit le motif d'échec. Ainsi, un attaquant ne peut pas deviner si un compte existe ou non dans le système. \\
\hline

Résolution transparente de l'identifiant & Connexion par email & Le candidat se connecte avec son email, mais le système traduit automatiquement cet email vers l'identifiant technique Salesforce, sans jamais exposer ce dernier. \\
\hline

QR~code éphémère & Connexion mobile & Les jetons QR ne sont valables que 5~minutes et sont chiffrés (SHA-256). La validation s'effectue côté serveur pour éviter toute falsification. \\
\hline

Journal d'audit & Toutes les opérations sensibles & Chaque connexion, réinitialisation ou tentative échouée est enregistrée avec horodatage, adresse email et niveau de sévérité (information, avertissement, critique). \\
\hline

Contrôle d'accès & Portail candidat & Deux jeux de permissions dédiés gouvernent les droits : l'un pour les candidats B2C, l'autre pour les objets de sécurité. Le visiteur non authentifié n'a accès qu'aux fonctionnalités publiques. \\
\hline

Validation côté serveur & Formulaire de candidature & Les types de fichier (PDF, DOC, DOCX, TXT) et la taille ($\leq$~10~Mo) sont vérifiés côté serveur, en complément de la validation côté navigateur, pour prévenir toute manipulation. \\
\hline

Isolation des données & Espace personnel & Chaque candidat ne peut consulter que ses propres données. L'isolation est garantie par les mécanismes natifs de partage de Salesforce (Sharing Sets). \\
\hline

\end{longtable}

\FloatBarrier

\section{Conclusion}

Ce premier sprint a permis de poser les fondations opérationnelles du \textbf{HR Talent Portal}. Le portail de recrutement est désormais accessible au public : un visiteur peut découvrir les offres d'emploi de Talan, les filtrer selon ses critères et soumettre sa candidature avec son CV. Parallèlement, l'infrastructure d'identité numérique est en place : le candidat peut se connecter par email ou par QR~code, réinitialiser son mot de passe en toute sécurité, et accéder à un tableau de bord personnel affichant le pipeline de son recrutement.

\medskip

Un soin particulier a été apporté à la sécurité : jetons cryptographiques à usage unique, messages d'erreur anti-énumération, journal d'audit traçant chaque opération sensible, et contrôle d'accès granulaire. Au total, ce sprint a produit \textbf{14~composants d'interface}, \textbf{5~contrôleurs serveur}, \textbf{3~entités de sécurité} et \textbf{2~jeux de permissions}.

\medskip

Ces fondations étant posées, le chapitre suivant sera consacré au Sprint~2 : l'intégration du moteur d'intelligence artificielle qui analysera automatiquement les CV, qualifiera les candidats et générera des descriptions d'offres professionnelles.

%%% Local Variables:
%%% mode: latex
%%% TeX-master: "isae-report-template"
%%% End:
